% does the Veritas system succeed in demonstrating the viability of a refinement type system with SMT-based verification for a sysems language toolchain?
% Experimental setup
% Program snippet examples
% Safety properties

\chapter{Evaluation}\label{ch:eval}

\section{Success criteria and methodology}

The aim of the evaluation is to test whether the implementation meets the core
goals stated in Chapter~\ref{ch:intro}: a type-preserving compilation pipeline
for a small systems language, an independent verifier for the low-level output,
and a practical enough implementation that the approach is not purely
theoretical. Following the plan set out in Chapter~\ref{ch:prep}, the
evaluation is organised around correctness, expressiveness, and overhead.

The evaluation uses four named suites plus a small robustness track:
\texttt{feature\_suite} (28 representative valid programs),
\texttt{negative\_suite} (35 intentionally invalid programs),
\texttt{verification\_stress\_suite} (6 proof-heavy programs), and
\texttt{algorithm\_suite} (4 runtime/practicality kernels). Chapter-labelled
runs were collected on 1 May 2026 using the manifest-driven harness in
\texttt{eval/bench.sh}. Compilation and verifier timings use built-in
\texttt{--bench} instrumentation, while end-to-end wall-clock timings use
\texttt{hyperfine} with 10 runs and 3 warmup runs. Runtime timings were
collected with \texttt{hyperfine}'s shell bypass enabled so the measurements
reflect the binaries rather than shell startup. Each run directory records the
git commit, tool versions, machine information, commands, frozen suite
manifests, and raw outputs.

\begin{table}[h]
  \centering
  \begin{tabular}{l l l}
    \hline
    \textbf{Criterion}                                          & \textbf{Evidence}         & \textbf{Current result} \\
    \hline
    Independent verifier agrees on supported fragment           & \texttt{feature\_suite}   & 28/28 both pass         \\
    Invalid programs are rejected                               & \texttt{negative\_suite}  & 35/35 rejected          \\
    Source language is expressive enough for non-trivial proofs & representative programs   & achieved                \\
    Verification overhead is measurable and attributable        & timing + SMT metrics      & achieved                \\
    Compiled kernels remain practical to run                    & \texttt{algorithm\_suite} & partial but positive    \\
    Implementation is robust on generated small programs        & seeded corpus             & 3/3 properties pass     \\
    \hline
  \end{tabular}
  \caption{Summary of the main evaluation criteria and their current status.}
  \label{tab:eval-summary}
\end{table}

\section{Functional correctness and trust architecture}

The most important correctness claim is local and explicit: for the frozen
\texttt{feature\_suite}, every representative source-language program accepted
by the frontend is also accepted by the independent physical DTAL verifier. The
chapter-labelled verifier-agreement run recorded 28 out of 28 programs as
\texttt{both\_pass}. This means that, within the currently supported fragment,
the trust architecture is functioning as intended: the compiler is untrusted,
but its output is independently re-checked before execution.

This feature suite is not a collection of trivial arithmetic programs. It
contains explicit conditionals and comparisons, arrays, references, mutable
state, preconditions, loop invariants, quantifiers, sortedness reasoning,
scalar dereference, safe division, overflow-aware arithmetic, and
borrow-passing examples. Table
~\ref{tab:expressive-programs} lists a representative subset of the programs
used in the chapter discussion.

\begin{table}[h]
  \centering
  \begin{tabular}{l l l}
    \hline
    \textbf{Program}                  & \textbf{Main feature(s)} & \textbf{Role in evaluation}   \\
    \hline
    \texttt{16\_array\_assignment}    & arrays, mutation, proofs & array-update reasoning        \\
    \texttt{19\_quantifier\_showcase} & quantifiers, logic       & quantified proof obligations  \\
    \texttt{20\_binary\_search}       & invariants, sortedness   & non-trivial search proof      \\
    \texttt{21\_safe\_division}       & preconditions, arithmetic safety & positive safety baseline \\
    \texttt{22\_bubble\_sort}         & mutation, sortedness     & end-to-end algorithmic proof  \\
    \texttt{30\_i64\_safe\_midpoint}  & refined integers, overflow safety & overflow-aware arithmetic \\
    \texttt{34\_shared\_scalar\_deref} & shared borrow, dereference & scalar reference semantics \\
    \texttt{33\_mutable\_borrow}      & exclusive borrow         & ownership discipline          \\
    \hline
  \end{tabular}
  \caption{Representative programs illustrating the supported source-language
    fragment.}
  \label{tab:expressive-programs}
\end{table}

The negative-suite results support the complementary claim that unsafe or
ill-typed programs are rejected systematically rather than on a few hand-picked
examples. All 35 programs in the frozen \texttt{negative\_suite} were
rejected. The suite spans 14 error classes, including type errors, name
resolution failures, array errors, contract violations, invariant failures, and
arithmetic/overflow safety errors. The breakdown is shown in Table
~\ref{tab:error-classes}.

\begin{table}[h]
  \centering
  \small
  \begin{tabular}{l r r}
    \hline
    \textbf{Error class}         & \textbf{Count} & \textbf{Rejected} \\
    \hline
    \texttt{type\_error}         & 7              & 7                 \\
    \texttt{name\_resolution}    & 2              & 2                 \\
    \texttt{mutability}          & 1              & 1                 \\
    \texttt{control\_flow}       & 3              & 3                 \\
    \texttt{array\_error}        & 2              & 2                 \\
    \texttt{call\_error}         & 2              & 2                 \\
    \texttt{contract\_violation} & 2              & 2                 \\
    \texttt{quantifier\_error}   & 3              & 3                 \\
    \texttt{invariant\_failure}  & 2              & 2                 \\
    \texttt{safety\_proof}       & 2              & 2                 \\
    \texttt{arithmetic\_safety}  & 2              & 2                 \\
    \texttt{overflow\_safety}    & 5              & 5                 \\
    \texttt{logic\_error}        & 1              & 1                 \\
    \texttt{mixed\_frontend}     & 1              & 1                 \\
    \hline
  \end{tabular}
  \caption{Negative-suite rejection results grouped by error class.}
  \label{tab:error-classes}
\end{table}

For runtime semantic agreement, all programs in the current runtime corpus
passed their expected semantic checks. Trivial examples were validated by exit
status where the exit code itself is the intended result, while algorithmic
kernels were checked by deterministic checksums. This is stronger evidence than
the earlier exit-only approach because it distinguishes successful execution
from accidental agreement on process status.

Finally, a separate fixed-seed robustness track was used to test
implementation-level stability rather than benchmark performance. The current
corpus contains 16 generated small programs across four families
(\texttt{valid\_arithmetic}, \texttt{valid\_arrays}, \texttt{invalid\_type},
\texttt{invalid\_syntax}). Three explicit robustness properties pass: valid
generated programs compile deterministically and verify in both in-memory and
text DTAL form, invalid generated programs fail deterministically with stable
diagnostics, and success/failure classification is stable across repeated
compilation. Notably, this track exposed a real determinism bug in synthesized
refinement naming, which was then fixed in the pipeline.

\section{Expressiveness}

The evaluation demonstrates that Veritas can express more than simple
first-order arithmetic properties. The supported fragment includes fixed-size
arrays, mutable updates, references and borrow-like alias restrictions,
preconditions, loop invariants, bounded quantifiers, specification-level
sortedness properties, safe integer reasoning, and scalar dereference through
shared and mutable borrows. Programs such as \texttt{20\_binary\_search},
\texttt{21\_safe\_division}, \texttt{22\_bubble\_sort},
\texttt{30\_i64\_safe\_midpoint}, and \texttt{34\_shared\_scalar\_deref} are
important here because they require multiple features to interact: arithmetic
facts must be preserved through loops, division must be guarded by
proof-carrying preconditions, array accesses must remain in bounds, borrows
must obey ownership discipline, and user-supplied invariants or quantified
specifications must be strong enough for verification to succeed.

This is also where the external kernels matter. The current external
runtime/practicality evidence is still modest, but the presence of translated
PolyBench-style kernels such as Floyd--Warshall and the two Jacobi benchmarks
shows that the system can handle structured numeric kernels with loop nests and
array updates, not just hand-crafted micro examples. That said, the evaluation
should not over-claim expressiveness beyond the currently frozen suites. The
language does not yet offer a complete story for arbitrary heap-heavy,
pointer-heavy, or concurrency-oriented systems programs, and those remain
outside the supported fragment.

\section{Verification and compilation overhead}

The feature-suite summary shows that the implementation is practical enough to
use, but that verification cost is highly non-uniform. Across the 28-program
feature suite, the mean source-to-output compile time was 25.870\,ms and the
mean explicit physical verification time was 19.625\,ms. The mean frontend SMT
time was 10.739\,ms and the mean verifier SMT time was 18.981\,ms. This already
indicates that solver work, rather than parsing or code emission, dominates the
cost profile on the interesting examples.

The most expensive programs make this clearer. Table~\ref{tab:hotspots}
highlights the main cost hotspots and the dominant reason for their expense.

\begin{table}[h]
  \centering
  \small
  \begin{tabular}{l r r l}
    \hline
    \textbf{Program}                  & \textbf{Compile (ms)} & \textbf{Verify (ms)} & \textbf{Dominant cost}    \\
    \hline
    \texttt{15\_array\_init}          & 40.654                & 292.555              & array proof obligations   \\
    \texttt{20\_binary\_search}       & 186.562               & 55.515               & invariants and sortedness \\
    \texttt{22\_bubble\_sort}         & 75.180                & 37.640               & invariants and sortedness \\
    \texttt{19\_quantifier\_showcase} & 36.185                & 24.616               & quantifiers               \\
    \texttt{20\_array\_loop\_invariant} & 63.534              & 21.393               & array proof obligations   \\
    \hline
  \end{tabular}
  \caption{Selected verification and compilation hotspots from the feature
    suite.}
  \label{tab:hotspots}
\end{table}

Two observations follow. First, the frontend and the independent verifier do
not incur the same cost on the same programs: \texttt{20\_binary\_search} is
the worst source-level compile case, while \texttt{15\_array\_init} is the
worst explicit verifier case. Second, cost scales more with proof structure
than with program size. Array initialisation and invariant-heavy algorithms are
much more expensive than straight-line code, simple comparisons, or basic
borrow/dereference examples, even when the latter add meaningful language
coverage.

End-to-end timing of the frozen \texttt{verification\_stress\_suite} reinforces
this point. Table~\ref{tab:timing-ratios} reports the ratio between
\texttt{compile\_plus\_verify} and \texttt{compile\_only} wall-clock timings.

\begin{table}[h]
  \centering
  \small
  \begin{tabular}{l r r r}
    \hline
    \textbf{Program}                  & \textbf{Compile only (s)} & \textbf{Compile+verify (s)} & \textbf{Ratio} \\
    \hline
    \texttt{15\_array\_init}          & 0.0648                    & 0.2280                      & 3.521          \\
    \texttt{16\_array\_assignment}    & 0.0890                    & 0.1089                      & 1.223          \\
    \texttt{19\_quantifier\_showcase} & 0.1607                    & 0.1835                      & 1.142          \\
    \texttt{20\_binary\_search}       & 0.2473                    & 0.2516                      & 1.017          \\
    \texttt{22\_bubble\_sort}         & 0.0848                    & 0.1175                      & 1.385          \\
    \texttt{sortedness}               & 0.0283                    & 0.0394                      & 1.393          \\
    \hline
  \end{tabular}
  \caption{End-to-end timing summary for the verification-stress suite.}
  \label{tab:timing-ratios}
\end{table}

The key result here is that verification overhead is not well-described by a
single multiplier. \texttt{15\_array\_init} is an outlier at 3.5$\times$,
whereas the remaining cases are much closer to parity. This supports the claim
that proof obligation structure, not merely the existence of verification, is
the main determinant of cost.

Runtime practicality was evaluated on four algorithm-suite kernels: a local
bubble-sort example plus Floyd--Warshall, Jacobi-1D, and Jacobi-2D. Table
~\ref{tab:runtime} shows the current runtime summary against GCC baselines.

\begin{table}[h]
  \centering
  \small
  \begin{tabular}{l r r r}
    \hline
    \textbf{Benchmark}             & \textbf{Veritas (s)} & \textbf{GCC \texttt{-O2} (s)} & \textbf{Ratio} \\
    \hline
    \texttt{bubble\_sort\_feature} & 0.000217             & 0.000590                      & 0.369          \\
    \texttt{floyd\_warshall\_64}   & 0.001878             & 0.001014                      & 1.853          \\
    \texttt{jacobi\_1d\_128\_80}   & 0.000374             & 0.000524                      & 0.713          \\
    \texttt{jacobi\_2d\_64\_20}    & 0.001584             & 0.000845                      & 1.875          \\
    \hline
  \end{tabular}
  \caption{Algorithm-suite runtime summary against GCC \texttt{-O2}.}
  \label{tab:runtime}
\end{table}

These numbers are best interpreted as supporting practicality evidence rather
than definitive performance claims. The larger kernels, especially
Floyd--Warshall and Jacobi-2D, show that verified generated binaries remain in
the same rough order of magnitude as compiled C baselines, albeit slower. The
smaller kernels are sub-millisecond and therefore much more sensitive to noise.
Hyperfine also reported outliers during some runtime measurements, so these
results should be read conservatively.

\section{Code size and trusted computing base}

Two distinct size questions matter here: how large the generated binaries are,
and how large the trusted computing base remains after moving verification to
physical DTAL.

Binary-size comparisons must be interpreted carefully. Total ELF size is not a
fully fair comparison because Veritas emits freestanding binaries, whereas GCC
links hosted startup and runtime code. For this reason, the \texttt{.text}
column is the more meaningful measure of code-generation overhead. Table
~\ref{tab:binary-size} summarises the current results.

\begin{table}[h]
  \centering
  \small
  \begin{tabular}{l r r r r}
    \hline
    \textbf{Program}             & \textbf{Veritas ELF} & \textbf{GCC \texttt{-O2} ELF} & \textbf{Veritas \texttt{.text}} & \textbf{GCC \texttt{-O2} \texttt{.text}} \\
    \hline
    \texttt{01\_simple}          & 623                  & 15720                         & 623                             & 265                                      \\
    \texttt{07\_function\_calls} & 881                  & 15816                         & 881                             & 311                                      \\
    \texttt{14\_for\_loops}      & 1339                 & 15712                         & 1339                            & 265                                      \\
    \texttt{22\_bubble\_sort}    & 1641                 & 15920                         & 1641                            & 441                                      \\
    \hline
  \end{tabular}
  \caption{Binary size summary. Total ELF size is not directly comparable
    across freestanding and hosted outputs; \texttt{.text} is the fairer code-size
    comparison.}
  \label{tab:binary-size}
\end{table}

The generated \texttt{.text} sections are larger than the GCC baselines for
these examples, which is consistent with Veritas carrying a more explicit
low-level proof structure and a simpler code-generation strategy. However, the
important trust-model result lies in the implementation size of the checker
rather than the size of any one output binary.

The current TCB summary is given in Table~\ref{tab:tcb}. The verifier occupies
7,864 lines out of 42,775 lines counted by the harness, i.e.\ 18.4\% of the
measured Rust codebase.

\begin{table}[h]
  \centering
  \begin{tabular}{l r r r}
    \hline
    \textbf{Trusted component} & \textbf{Lines} & \textbf{Total lines} & \textbf{Trusted ratio} \\
    \hline
    \texttt{verifier\_tcb}     & 7864           & 42775                & 0.184                  \\
    \hline
  \end{tabular}
  \caption{Current trusted computing base summary from the evaluation harness.}
  \label{tab:tcb}
\end{table}

This supports the central architectural claim of the project: the trusted core
is much smaller than the overall compiler implementation, and the untrusted
parts of the backend can be pushed below the verification boundary.

\section{Limitations and current issues}

Several limitations remain and should be stated plainly.

First, the evaluation demonstrates structural safety, not full semantic
compiler correctness. The independent verifier can reject malformed low-level
output, but this is not the same as a proof that compilation preserves the
entire meaning of every source program.

Second, the runtime benchmark set is still small. The current external kernels
cover only the committed translated workloads in the repository. This is enough
to support a modest practicality claim, but not enough to claim broad
competitiveness against mature systems toolchains.

Third, not all quantitative comparisons are equally strong. Runtime results for
sub-millisecond kernels are noisy, and total ELF size comparisons are affected
by freestanding versus hosted startup/runtime differences. These are useful
measurements, but they require caveats in interpretation.

Fourth, the robustness track is deliberately narrow. The fixed-seed generated
corpus is useful for determinism and crash-resistance testing, but it does not
cover the full source-language fragment, external benchmarks, or arbitrary DTAL
inputs. It strengthens confidence in the implementation; it does not replace
the benchmark suites.

Finally, the current supported fragment is still intentionally limited. The
project does not yet offer a complete story for arbitrary heap-heavy systems
programs, concurrency properties, or a fully general low-level memory model.
Those remain natural directions for future work.

\section{Summary against objectives}

Against the core aims stated in Chapter~\ref{ch:intro}, the project can be
judged a success with clear boundaries.

The first objective was to build a type-preserving compiler pipeline for a
small systems language. This has been achieved: the source language supports a
non-trivial fragment with arrays, mutable state, contracts, loop invariants,
quantifiers, and borrow-like features, and these properties are carried through
to DTAL.

The second objective was to build an independent verifier that checks the
compiler's output without access to the source program. This has also been
achieved for the supported fragment: all 22 programs in the frozen feature
suite were accepted by both the frontend and the physical DTAL verifier, and
the negative suite was rejected systematically.

The third objective was to evaluate the practicality and limits of the
architecture. The results here are positive but more qualified. Verification is
practical on the current suite, but its cost depends strongly on proof
structure. Compiled kernels remain runnable and in the same broad performance
region as simple C baselines, but the benchmark set is still too small and too
noisy in places for stronger performance claims.

Overall, the evaluation supports the central thesis of the dissertation: it is
possible to move trust away from a large compiler and into a smaller
independent verifier while retaining useful source-language expressiveness and
without making the resulting toolchain prohibitively expensive to use.
